\documentclass[a4paper,11pt]{article}
\usepackage{graphicx}
\input{../outputs/results.tex}
\title{Example Loss Calculation}
\author{LossCalculation-SPB}
\date{\today}
\begin{document}
\maketitle
\section{Model and assumptions}
For a constant resistance, the average conduction loss is
\begin{equation}
P_{\mathrm{loss}}=I_{\mathrm{RMS}}^2 R.
\end{equation}
This illustrative sweep assumes $R=\Resistance\,\Omega$ and a current
between zero and $\MaxCurrent\,\mathrm{A}$ RMS. Temperature dependence,
switching losses and other parasitic effects are excluded.
\section{Results}
At $I_{\mathrm{RMS}}=\MaxCurrent\,\mathrm{A}$, the calculated loss is
$\MaxLoss\,\mathrm{W}$. The values and figure below are generated by Python.
\begin{figure}[ht]
\centering
\includegraphics[width=0.9\linewidth]{../outputs/loss_plot.pdf}
\caption{Conduction loss as a function of RMS current.}
\end{figure}
\end{document}
